% trimer_setup.tex — Chapter 4: Trimer BSSE Analysis — Setup and Geometry
% Project: rmrresearch/bsse_db | Branch: fcuantico
% Molecules: H2O (complete), Ne (complete), HF (complete)

% ======================================================================
\section{\ce{H2O} Trimer}
% ======================================================================

% ----------------------------------------------------------------------
\subsection{Overview}
% ----------------------------------------------------------------------

The \ce{H2O} trimer extends the dimer analysis to a three-body system.
Five O--O distances are studied: $R = 2.75$, $3.0$, $3.25$, $3.5$,
and \SI{3.75}{\angstrom}. All calculations were performed at the SCF,
MP2, and CCSD(T) levels using the aug-cc-pVDZ basis
set~\cite{dunning1989,kendall1992} with the NWChem quantum chemistry
package~\cite{nwchem2020}. The many-body BSSE decomposition follows
the VMFC scheme~\cite{valiron1997}.

% ----------------------------------------------------------------------
\subsection{Geometry}
% ----------------------------------------------------------------------

The three \ce{H2O} monomers, labeled A, B, and C, are arranged in a
planar equilateral triangle with all atoms in the $xy$ plane ($z = 0$).
Within each monomer the O--H bond lengths and H--O--H angle are held
fixed at their equilibrium values, identical to those used in the dimer
calculations. Table~\ref{tab:h2o_trimer_geometry} lists the Cartesian
coordinates at $R = \SI{2.75}{\angstrom}$, the shortest O--O distance
in the dataset. For the remaining configurations the triangle is
uniformly expanded while preserving the equilateral arrangement and
the monomer internal geometry.

\begin{table}[!ht]
    \centering
    \caption{Cartesian coordinates of the \ce{H2O} trimer at
    $R = \SI{2.75}{\angstrom}$ (aug-cc-pVDZ, NWChem convention).
    Atom labels: A = monomer~A (oxygen at origin), B = monomer~B
    (oxygen at upper vertex), C = monomer~C (oxygen at right vertex).
    All coordinates in \AA.}
    \label{tab:h2o_trimer_geometry}
    \begin{tabular}{lcccc}
        \toprule
        Atom & Monomer & $x$ (\AA) & $y$ (\AA) & $z$ (\AA) \\
        \midrule
        O & A & $-0.15219$ & $-0.00116$ & $0.00000$ \\
        H & A & $ 0.43777$ & $-0.76611$ & $0.00000$ \\
        H & A & $ 0.44760$ & $ 0.75608$ & $0.00000$ \\
        O & B & $ 1.47281$ & $ 2.81342$ & $0.00000$ \\
        H & B & $ 2.06277$ & $ 2.04847$ & $0.00000$ \\
        H & B & $ 2.07260$ & $ 3.57066$ & $0.00000$ \\
        O & C & $ 3.09781$ & $-0.00116$ & $0.00000$ \\
        H & C & $ 3.68777$ & $-0.76611$ & $0.00000$ \\
        H & C & $ 3.69760$ & $ 0.75608$ & $0.00000$ \\
        \bottomrule
    \end{tabular}
\end{table}



% TODO: move this table to theory.tex when that chapter is written
\subsection{VMFC File Keys}
% ----------------------------------------------------------------------

The VMFC scheme for a three-monomer system (A, B, C) requires a total
of 19 distinct NWChem calculations per configuration. Each calculation
is identified by a \textit{file key} of the form \texttt{real\_basis},
where \texttt{real} denotes the monomer(s) treated as real atoms and
\texttt{basis} denotes the set of monomers contributing basis
functions. Table~\ref{tab:h2o_trimer_filekeys} lists all 19 file keys
and their physical meaning.

\begin{table}[!ht]
    \centering
    \caption{The 19 VMFC file keys for the \ce{H2O} trimer. The
    \texttt{real\_basis} naming convention follows the folder structure
    in \texttt{bsse\_db}: \texttt{real} = real atoms, \texttt{basis} =
    basis set used (own + ghost atoms). Keys with identical monomers
    (e.g.\ \texttt{B\_B} and \texttt{A\_A}) are verified numerically
    to confirm monomer equivalence.}
    \label{tab:h2o_trimer_filekeys}
    \small
    \begin{tabular}{lll}
        \toprule
        File key & Real atoms & Basis set \\
        \midrule
        \texttt{A\_A}    & A     & A only \\
        \texttt{B\_B}    & B     & B only \\
        \texttt{C\_C}    & C     & C only \\
        \texttt{AB\_AB}  & A, B  & AB \\
        \texttt{AC\_AC}  & A, C  & AC \\
        \texttt{BC\_BC}  & B, C  & BC \\
        \texttt{ABC\_ABC}& A, B, C & ABC (full trimer) \\
        \texttt{A\_AB}   & A     & AB (B as ghost) \\
        \texttt{B\_AB}   & B     & AB (A as ghost) \\
        \texttt{A\_AC}   & A     & AC (C as ghost) \\
        \texttt{C\_AC}   & C     & AC (A as ghost) \\
        \texttt{B\_BC}   & B     & BC (C as ghost) \\
        \texttt{C\_BC}   & C     & BC (B as ghost) \\
        \texttt{A\_ABC}  & A     & ABC (B, C as ghosts) \\
        \texttt{B\_ABC}  & B     & ABC (A, C as ghosts) \\
        \texttt{C\_ABC}  & C     & ABC (A, B as ghosts) \\
        \texttt{AB\_ABC} & A, B  & ABC (C as ghost) \\
        \texttt{AC\_ABC} & A, C  & ABC (B as ghost) \\
        \texttt{BC\_ABC} & B, C  & ABC (A as ghost) \\
        \bottomrule
    \end{tabular}
\end{table}

The 19 file keys divide naturally into four groups by the size of the
basis set used: the three monomer-in-own-basis calculations
(\texttt{A\_A}, \texttt{B\_B}, \texttt{C\_C}), the six
monomer-in-dimer-basis calculations (\texttt{A\_AB}, \texttt{B\_AB},
etc.), the three monomer-in-trimer-basis calculations
(\texttt{A\_ABC}, \texttt{B\_ABC}, \texttt{C\_ABC}), and the three
dimer-in-trimer-basis calculations (\texttt{AB\_ABC}, \texttt{AC\_ABC},
\texttt{BC\_ABC}). Together these calculations provide all the
ingredients needed to evaluate both FORM~I and FORM~II of the VMFC
BSSE decomposition.

% ----------------------------------------------------------------------
\subsection{Configuration Summary}
% ----------------------------------------------------------------------

Table~\ref{tab:h2o_trimer_configs} summarizes the five \ce{H2O} trimer
configurations studied in this work.

\begin{table}[!ht]
    \centering
    \caption{Summary of the five \ce{H2O} trimer configurations.
    All configurations use the aug-cc-pVDZ basis set and were computed
    at the SCF, MP2, and CCSD(T) levels of theory.}
    \label{tab:h2o_trimer_configs}
    \begin{tabular}{ccc}
        \toprule
        Configuration & O--O Distance (\AA) & File key count \\
        \midrule
        0 & 2.75 & 19 \\
        1 & 3.00 & 19 \\
        2 & 3.25 & 19 \\
        3 & 3.50 & 19 \\
        4 & 3.75 & 19 \\
        \bottomrule
    \end{tabular}
\end{table}

% ======================================================================
\section{Ne Trimer}
% ======================================================================

% ----------------------------------------------------------------------
\subsection{Overview}
% ----------------------------------------------------------------------

The Ne trimer extends the dimer analysis to a three-body system of
rare-gas atoms. Ten Ne--Ne distances are studied: $R = 1.5$, $1.75$,
$2.0$, $2.25$, $2.5$, $2.75$, $3.0$, $3.25$, $3.5$, and
\SI{3.75}{\angstrom}. All calculations were performed at the SCF, MP2,
and CCSD(T) levels using the aug-cc-pVDZ basis
set~\cite{dunning1989,kendall1992} with the NWChem quantum chemistry
package~\cite{nwchem2020}. The many-body BSSE decomposition follows
the VMFC scheme~\cite{valiron1997}.

% ----------------------------------------------------------------------
\subsection{Geometry}
% ----------------------------------------------------------------------

The three Ne atoms, labeled A, B, and C, are arranged in a planar
equilateral triangle with all atoms in the $xy$ plane ($z = 0$).
Table~\ref{tab:ne_trimer_geometry} lists the Cartesian coordinates at
$R = \SI{1.5}{\angstrom}$, the shortest Ne--Ne distance in the dataset.
For the remaining configurations the triangle is uniformly expanded
while preserving the equilateral arrangement. The three pairwise
distances AB, AC, and BC are identical by construction at every
configuration.

\begin{table}[!ht]
    \centering
    \caption{Cartesian coordinates of the Ne trimer at
    $R = \SI{1.5}{\angstrom}$ (aug-cc-pVDZ, NWChem convention).
    Atom labels: A = monomer~A (at origin), B = monomer~B
    (upper vertex), C = monomer~C (right vertex).
    All coordinates in \AA.}
    \label{tab:ne_trimer_geometry}
    \begin{tabular}{lcccc}
        \toprule
        Atom & Monomer & $x$ (\AA) & $y$ (\AA) & $z$ (\AA) \\
        \midrule
        Ne & A & $0.00000000$ & $0.00000000$ & $0.00000000$ \\
        Ne & B & $0.75000000$ & $1.29903811$ & $0.00000000$ \\
        Ne & C & $1.50000000$ & $0.00000000$ & $0.00000000$ \\
        \bottomrule
    \end{tabular}
\end{table}

Because all three Ne monomers are identical, the three monomer
energies are numerically equal: $E_\text{A}(\text{A}) =
E_\text{B}(\text{B}) = E_\text{C}(\text{C})$ at every level of theory
and every configuration. This identity is verified numerically in the
pipeline and used to confirm the integrity of the parsed data.

% ----------------------------------------------------------------------
\subsection{Configuration Summary}
% ----------------------------------------------------------------------

Table~\ref{tab:ne_trimer_configs} summarizes the ten Ne trimer
configurations studied in this work.

\begin{table}[!ht]
    \centering
    \caption{Summary of the ten Ne trimer configurations.
    All configurations use the aug-cc-pVDZ basis set and were computed
    at the SCF, MP2, and CCSD(T) levels of theory.}
    \label{tab:ne_trimer_configs}
    \begin{tabular}{ccc}
        \toprule
        Configuration & Ne--Ne Distance (\AA) & File key count \\
        \midrule
        0 & 1.50 & 19 \\
        1 & 1.75 & 19 \\
        2 & 2.00 & 19 \\
        3 & 2.25 & 19 \\
        4 & 2.50 & 19 \\
        5 & 2.75 & 19 \\
        6 & 3.00 & 19 \\
        7 & 3.25 & 19 \\
        8 & 3.50 & 19 \\
        9 & 3.75 & 19 \\
        \bottomrule
    \end{tabular}
\end{table}

% ======================================================================
\section{HF Trimer}
% ======================================================================

% ----------------------------------------------------------------------
\subsection{Overview}
% ----------------------------------------------------------------------

The HF trimer extends the dimer analysis to a three-body system of
hydrogen fluoride molecules. Seven F--F distances are studied:
$R = 1.5$, $1.75$, $2.0$, $2.25$, $2.5$, $2.75$, and
\SI{3.0}{\angstrom}. All calculations were performed at the SCF, MP2,
and CCSD(T) levels using the aug-cc-pVDZ basis
set~\cite{dunning1989,kendall1992} with the NWChem quantum chemistry
package~\cite{nwchem2020}. The many-body BSSE decomposition follows
the VMFC scheme~\cite{valiron1997}.

% ----------------------------------------------------------------------
\subsection{Geometry}
% ----------------------------------------------------------------------

The three HF monomers, labeled A, B, and C, are arranged so that the
three fluorine atoms form a planar equilateral triangle in the $xy$
plane ($z = 0$), with the hydrogen atom of each monomer displaced
along the negative $z$ direction. Within each monomer the H--F bond
length is held fixed. Table~\ref{tab:hf_trimer_geometry} lists the
Cartesian coordinates at $R = \SI{1.5}{\angstrom}$, the shortest F--F
distance in the dataset. For the remaining configurations the
fluorine triangle is uniformly expanded while preserving the
equilateral arrangement and the monomer internal geometry.

\begin{table}[!ht]
    \centering
    \caption{Cartesian coordinates of the HF trimer at
    $R = \SI{1.5}{\angstrom}$ (aug-cc-pVDZ, NWChem convention).
    Atom labels: A = monomer~A (F at origin), B = monomer~B
    (F at upper vertex), C = monomer~C (F at right vertex).
    The three F atoms lie in the $z = 0$ plane; the H atoms are
    displaced to $z = -0.924$~\AA. All coordinates in \AA.}
    \label{tab:hf_trimer_geometry}
    \begin{tabular}{lcccc}
        \toprule
        Atom & Monomer & $x$ (\AA) & $y$ (\AA) & $z$ (\AA) \\
        \midrule
        F & A & $0.00000000$ & $0.00000000$ & $ 0.00000000$ \\
        H & A & $0.00000000$ & $0.00000000$ & $-0.92400000$ \\
        F & B & $0.75000000$ & $1.29903811$ & $ 0.00000000$ \\
        H & B & $0.75000000$ & $1.29903811$ & $-0.92400000$ \\
        F & C & $1.50000000$ & $0.00000000$ & $ 0.00000000$ \\
        H & C & $1.50000000$ & $0.00000000$ & $-0.92400000$ \\
        \bottomrule
    \end{tabular}
\end{table}

Although the HF monomer is a diatomic and the three monomers are
individually equivalent, the H atoms point out of the fluorine plane,
breaking the in-plane $C_{3v}$ symmetry present in the Ne trimer.
Nevertheless, the three F--F pairwise distances AB, AC, and BC remain
identical by construction at every configuration, and the three
monomer energies are numerically equal:
$E_\text{A}(\text{A}) = E_\text{B}(\text{B}) = E_\text{C}(\text{C})$
at every level of theory.

% ----------------------------------------------------------------------
\subsection{Configuration Summary}
% ----------------------------------------------------------------------

Table~\ref{tab:hf_trimer_configs} summarizes the seven HF trimer
configurations studied in this work.

\begin{table}[!ht]
    \centering
    \caption{Summary of the seven HF trimer configurations.
    All configurations use the aug-cc-pVDZ basis set and were computed
    at the SCF, MP2, and CCSD(T) levels of theory.}
    \label{tab:hf_trimer_configs}
    \begin{tabular}{ccc}
        \toprule
        Configuration & F--F Distance (\AA) & File key count \\
        \midrule
        0 & 1.50 & 19 \\
        1 & 1.75 & 19 \\
        2 & 2.00 & 19 \\
        3 & 2.25 & 19 \\
        4 & 2.50 & 19 \\
        5 & 2.75 & 19 \\
        6 & 3.00 & 19 \\
        \bottomrule
    \end{tabular}
\end{table}